% Tesina de Licenciatura — Facultad de Informática, UNLP
% Compilar: latexmk -pdf -output-directory=build main.tex
\documentclass[12pt,a4paper,oneside]{book}

\usepackage[utf8]{inputenc}
\usepackage[T1]{fontenc}
\usepackage[spanish,es-tabla]{babel}
\usepackage[left=3cm,right=2.5cm,top=2.5cm,bottom=2.5cm]{geometry}
\usepackage{graphicx}
\usepackage{booktabs}
\usepackage{amsmath}
\usepackage[hidelinks]{hyperref}
\usepackage{csquotes}

\title{Modelo vs.\ contexto: estudio comparativo del impacto del LLM y del
  conocimiento de dominio en la generación de código por agentes de IA}
\author{Mariano Alex Martinelli}
\date{\today}

\begin{document}

\frontmatter
\maketitle
% \input{capitulos/00-agradecimientos}
\tableofcontents

\mainmatter
\chapter{Introducción}\label{cap:introduccion}

% Fuente: sección "Motivación" de la propuesta (Formulario 1).

\section{Motivación}
% Del desarrollo tradicional a los agentes de IA; el vacío empírico: dado un pipeline
% elegido, ¿cuánto impactan modelo y conocimiento de dominio en la calidad del código?

\section{Pregunta de investigación y objetivos}
% Objetivo principal (factorial 2×2) + objetivos específicos de la propuesta.

\section{Contribuciones esperadas}
% Spec/holdout reutilizable, dataset comparativo, análisis de efectos, recomendaciones.

\section{Estructura del documento}

\chapter{Estado del arte y marco teórico}\label{cap:estado-del-arte}

% Redactable desde ya, sin dependencia de las corridas.

\section{LLMs para generación de código}
% Chen et al. 2021; benchmarks (HumanEval, MBPP, SoftwareDev) y sus límites.

\section{Agentes de IA para ingeniería de software}
% MetaGPT, ChatDev, SWE-Agent, AgentMesh; survey Wang et al. 2024.

\section{Frameworks de agentes}
% Claude Agent SDK, OpenAI Agents SDK.

\section{Retrieval-Augmented Generation}
% Lewis et al. 2020; RAG como inyección de conocimiento de dominio.

\section{Exchanges centralizados de criptomonedas}
% Arquitectura de referencia; Harris 2003 (microestructura de mercado).

\section{Estándares on-chain: BIPs y EIPs}
% BIP-32/39/44, EIP-155, ERC-20; por qué son conocimiento discreto y verificable.

\section{El vacío empírico que aborda este trabajo}

\chapter{Caso de estudio: un exchange con capa on-chain mínima}\label{cap:caso-de-estudio}

% Depende de H1 (spec freeze) para citar la versión final de la spec.

\section{Alcance y justificación del dominio}
% Por qué un exchange; dentro/fuera de alcance; par único ETH/USDC-mock, Sepolia.

\section{Arquitectura de referencia}
% Matching engine, reglas de negocio, capa on-chain, clientes web/mobile.

\section{La especificación funcional y su doble rol}
% Input común de las corridas + holdout de evaluación; épicas y HUs; AT-ids estables.

\section{Convenciones normativas}
% Dinero en enteros, redondeo determinista, modelo de errores, invariantes INV-1..8.

\chapter{Diseño experimental}\label{cap:diseno-experimental}

% Depende de H2 (protocolo pre-registrado). Fuente: "Marco metodológico" de la propuesta.

\section{Marco metodológico}
% Design Science Research (Hevner et al. 2004) + estudio de caso (Yin 2018).

\section{Diseño factorial 2×2}
% Factores: modelo (A/B) × RAG (sin/con); las cuatro celdas.

\section{Variables dependientes}
% Tasa de ATs superados (holdout), intervenciones por causa raíz, alucinaciones de
% dominio, métricas estáticas, adherencia a estándares; rol desarrollador vs. revisor.

\section{Variables controladas y protocolo de corrida}
% Spec pinneada, corpus congelado, evaluador único, criterios de intervención
% pre-registrados, presupuestos, orden de construcción.

\section{Amenazas a la validez}
% n=1 por celda (sin inferencia estadística), deriva de modelos comerciales, sesgo
% del evaluador (mitigado por harness pre-construido), generalización del caso único.

\chapter{Infraestructura del experimento}\label{cap:infraestructura}

% Depende de H3–H5.

\section{Corpus de conocimiento de dominio}
% Curaduría de BIPs/EIPs, manifest, congelamiento.

\section{Harness de agentes}
% Claude Agent SDK y OpenAI Agents SDK; paridad entre condiciones; RAG conmutable.

\section{Harness de evaluación}
% Suite black-box por AT-id contra el contrato HTTP/WS; rúbricas web/mobile;
% detección de alucinaciones; métricas estáticas.

\section{Registro y trazabilidad}
% Manifests por corrida, log de intervenciones, journal y ADRs; repos por corrida.

\chapter{Resultados}\label{cap:resultados}

% Depende de H7–H8.

\section{Resumen de las corridas}
% Tabla por celda: duración, costo, tokens, commits, intervenciones.

\section{Tasa de tests de aceptación superados}
% Por celda, desagregada por épica/componente.

\section{Intervenciones humanas}
% Cantidad y distribución por causa raíz, por celda.

\section{Alucinaciones de dominio}

\section{Métricas estáticas de calidad}

\section{Rol revisor (análisis cualitativo)}

\chapter{Análisis y discusión}\label{cap:analisis}

% Depende de H9.

\section{Efecto principal del modelo}

\section{Efecto principal del RAG}

\section{Interacción modelo × RAG}

\section{Patrones cualitativos de fallo}
% Por categoría de intervención y por épica; casos ilustrativos.

\section{Implicancias prácticas}
% Cuándo invertir en inyección de conocimiento de dominio vs. elección del modelo.

\section{Limitaciones}

\chapter{Conclusiones y trabajo futuro}\label{cap:conclusiones}

\section{Conclusiones}

\section{Trabajo futuro}
% Réplicas por celda, más modelos, otros dominios con estándares formales,
% granularidad del corpus RAG, evaluación del rol revisor en profundidad.


\backmatter
\bibliographystyle{apalike}
\bibliography{bibliografia}
% Apéndices previstos: catálogo de AT-ids, manifests de corridas, protocolo experimental.

\end{document}
